\section{Experimental Methodology}

\subsection{Evaluation philosophy}

The campaign was designed to test whether conclusions drawn from a narrow clean-speech benchmark
survive changes in speech difficulty, recognizer, repeated inference, platform, and downstream
task.

We separate stable benchmark measurements from exploratory mechanism probes. Stable artifacts
provide the primary manuscript results. Follow-up probes are used when an unexpected result
requires an intervention to identify its mechanism. Derived significance files analyze existing
measurements and are not counted as independent experiments.

A central reporting rule is that every headline result must specify the population or corpus,
condition, uncertainty where available, and relevant caveat. Point-estimate rankings are not
interpreted as statistical superiority unless the corresponding paired or interval evidence
supports that interpretation.

\subsection{Speech-recognition corpora}

The principal ASR experiments use deterministic subsets of LibriSpeech
\citep{panayotov2015librispeech}. The clean condition uses 200 utterances from
\texttt{test-clean}, selected through the benchmark's deterministic speaker-stratified
procedure. The harder condition uses 200 utterances from \texttt{test-other}, covering 33
speakers and approximately 23.4 minutes of audio.

\texttt{test-other} is more difficult than \texttt{test-clean}, but it remains read audiobook
speech. We therefore use it as a robustness stress condition and do not interpret it as a direct
estimate of spontaneous desktop-dictation accuracy.

The benchmark records a digest of the selected utterance identifiers so that two artifacts
claiming to evaluate the same number of utterances can be checked for corpus identity rather than
assumed to have used identical input.

\subsection{Meeting corpora}

Real-meeting diarization is evaluated on the complete 16-recording AMI test split used by the
campaign. The set contains approximately 543.7 minutes of four-person meeting recordings, with
30,714 seconds of scored reference speech in the archived aggregate. Human RTTM annotations
provide the reference speaker labels.

The AMI condition used here is more representative of real conversational turn-taking than the
project's synthetic meeting fixtures, but the evaluated headset mix is still cleaner than a
typical single far-field laptop microphone. We therefore report AMI as real annotated meeting
evidence without treating it as a complete far-field deployment estimate.

Synthetic meetings and selected VoxConverse recordings serve different purposes. Synthetic audio
is retained as a controlled regression fixture. VoxConverse is used primarily to examine whether
clustering heuristics calibrated on meeting-like conditions transfer to more crowded speaker
distributions. Results from these domains are not pooled into one \der{}.

\subsection{Hardware and software conditions}

The main August engine-comparison campaign ran on a CPU-only Azure
\texttt{Standard\_D16s\_v6} host with 16 logical Intel Xeon Platinum 8573C CPUs and
approximately 67 GB of memory under Ubuntu 24.04. Inference used integer-8 execution where
supported. Exact library and \YazSes{} versions are preserved per result artifact rather than
assumed from one campaign-wide environment.

Cross-platform decoding was separately evaluated on Linux x86-64, Linux arm64, macOS arm64, and
Windows x86-64 runners using the same 60-utterance subset and CTranslate2 4.8.1 for the compared
faster-whisper checkpoints.

Timing results are interpreted within a machine condition. We do not use hosted-runner latency to
make direct absolute performance comparisons across heterogeneous hosts.

\subsection{Metrics}

Speech-recognition accuracy is measured with word error rate:
\begin{equation}
  \wer = \frac{S + D + I}{N},
\end{equation}
where $S$, $D$, and $I$ are substitution, deletion, and insertion counts and $N$ is the number
of reference words.

Because the repeated-decode investigation is specifically concerned with catastrophic extra text,
we retain the individual error components rather than relying on \wer{} alone.

Decode performance is reported using real-time factor. Lower values correspond to faster
transcription, but \rtf{} values are hardware-specific.

Paired speech-recognition comparisons use paired analysis when the same utterances are decoded
under two conditions. The onset experiment uses paired first-word outcomes and a
multiplicity-aware analysis rather than inferring onset effects from whole-utterance \wer{}.

Diarization is reported using \der{}, with both per-recording mean and time-weighted corpus
aggregates where appropriate. We separately retain speaker-count error because a system can
produce the right number of clusters while still assigning speech to the wrong people.

Streaming evaluation reports final speech-end-to-text latency, time to first useful partial
output, and the fraction of final text already visible at release.

\subsection{Reproducibility and provenance}

Every stable result is intended to answer both the measured value and the conditions that
produced it. Artifacts therefore retain the benchmark command, relevant software versions,
machine information, and corpus identity. The result archive also preserves displaced outputs
when a benchmark is repeated with the same nominal filename.

This mechanism matters scientifically. During the hard-speech engine campaign, a repeated
\texttt{large-v3} measurement disagreed substantially with an earlier run. Treating the newest
artifact as authoritative would have erased the disagreement. Retaining both instead converted
the contradiction into a reproducibility experiment.

The analysis below therefore treats disagreement between repeated runs as evidence to be
explained, not noise to be silently replaced.
